% !TEX program = pdflatex
%=====================================================================
%  Criptografía y Seguridad - ITBA
%  Clase 9 — Flujo de Información y Malware (Seguridad en aplicaciones)
%  Práctica: CWE Top 25 / errores de software más peligrosos
%  Apunte integral: teoría + práctica
%=====================================================================
\documentclass[11pt,a4paper]{article}

%---------------------------------------------------------------------
%  Paquetes
%---------------------------------------------------------------------
\usepackage[utf8]{inputenc}
\usepackage[T1]{fontenc}
\usepackage[spanish,es-tabla,es-noshorthands]{babel}
\usepackage{lmodern}
\usepackage{microtype}

\usepackage{amsmath,amssymb,amsthm}
\usepackage{mathtools}
\usepackage{geometry}
\geometry{a4paper,margin=2.4cm,headheight=22pt}

\usepackage{xcolor}
\usepackage{graphicx}
\usepackage{enumitem}
\usepackage{booktabs}
\usepackage{array}
\usepackage{tcolorbox}
\tcbuselibrary{skins,breakable}
\usepackage{fancyhdr}
\usepackage{titlesec}
\usepackage{hyperref}
\usepackage{listings}
\usepackage{caption}

\usepackage{tikz}
\usetikzlibrary{shapes.geometric,shapes.misc,shapes.symbols,arrows.meta,positioning,calc,
                fit,backgrounds,decorations.pathmorphing,shadows.blur,chains}

%---------------------------------------------------------------------
%  Paleta de color (inspirada en las diapositivas de la cátedra)
%---------------------------------------------------------------------
\definecolor{itbaCyan}{HTML}{6EC6D9}
\definecolor{itbaCyanDark}{HTML}{2E8B9E}
\definecolor{itbaBlue}{HTML}{AEDCEA}
\definecolor{boxBlue}{HTML}{BBD4F0}
\definecolor{slideGray}{HTML}{ECECEC}
\definecolor{practGreen}{HTML}{4F8A3D}
\definecolor{practGreenL}{HTML}{E2EFD9}
\definecolor{attackRed}{HTML}{C0392B}
\definecolor{codeBg}{HTML}{F5F5F5}
% colores de las categorías CWE (recreación de las slides prácticas)
\definecolor{cweNeutral}{HTML}{F6C85A}
\definecolor{cweResource}{HTML}{B79CD6}
\definecolor{cweAccess}{HTML}{8CC07A}
\definecolor{cweProtect}{HTML}{F0A868}
\definecolor{cweOther}{HTML}{D9D9D9}

%---------------------------------------------------------------------
%  Hipervínculos
%---------------------------------------------------------------------
\hypersetup{
  colorlinks=true,
  linkcolor=itbaCyanDark,
  urlcolor=itbaCyanDark,
  citecolor=itbaCyanDark,
  pdftitle={Clase 9 - Flujo de informacion y malware},
  pdfauthor={Criptografia y Seguridad - ITBA}
}

%---------------------------------------------------------------------
%  Encabezados y pies
%---------------------------------------------------------------------
\pagestyle{fancy}
\fancyhf{}
\lhead{\small\textsc{Criptografía y Seguridad — ITBA}}
\rhead{\small Clase 9 · Flujo de Información y Malware}
\cfoot{\small\thepage}
\renewcommand{\headrulewidth}{0.4pt}
\renewcommand{\headrule}{\hbox to\headwidth{\color{itbaCyanDark}\leaders\hrule height \headrulewidth\hfill}}

%---------------------------------------------------------------------
%  Títulos de sección con barra de color (estilo diapositiva)
%---------------------------------------------------------------------
\titleformat{\section}
  {\normalfont\Large\bfseries\color{itbaCyanDark}}
  {\thesection}{0.6em}{}
  [\vspace{-0.4em}{\color{itbaCyan}\titlerule[2pt]}]
\titleformat{\subsection}
  {\normalfont\large\bfseries\color{itbaCyanDark}}
  {\thesubsection}{0.6em}{}
\titleformat{\subsubsection}
  {\normalfont\normalsize\bfseries\color{black}}
  {\thesubsubsection}{0.6em}{}

%---------------------------------------------------------------------
%  Cajas reutilizables
%---------------------------------------------------------------------
% Caja "diapositiva" (recuadro gris de las slides)
\newtcolorbox{slidebox}[1][]{
  enhanced, colback=slideGray, colframe=black!55, boxrule=0.5pt,
  arc=1pt, left=6pt,right=6pt,top=4pt,bottom=4pt, #1}

% Caja de definición
\newtcolorbox{defbox}[1]{
  enhanced, breakable, colback=itbaBlue!25, colframe=itbaCyanDark,
  boxrule=1pt, arc=2pt, left=8pt,right=8pt,top=6pt,bottom=6pt,
  fonttitle=\bfseries, coltitle=white, title={#1}}

% Caja de nota práctica
\newtcolorbox{practbox}[1]{
  enhanced, breakable, colback=practGreenL, colframe=practGreen,
  boxrule=1pt, arc=2pt, left=8pt,right=8pt,top=6pt,bottom=6pt,
  fonttitle=\bfseries, coltitle=white, title={#1}}

% Caja de atención / ataque
\newtcolorbox{warnbox}[1]{
  enhanced, breakable, colback=attackRed!8, colframe=attackRed,
  boxrule=1pt, arc=2pt, left=8pt,right=8pt,top=6pt,bottom=6pt,
  fonttitle=\bfseries, coltitle=white, title={#1}}

%---------------------------------------------------------------------
%  Estilo de código
%---------------------------------------------------------------------
\lstdefinestyle{code}{
  basicstyle=\ttfamily\footnotesize,
  backgroundcolor=\color{codeBg},
  keywordstyle=\color{itbaCyanDark}\bfseries,
  commentstyle=\color{practGreen}\itshape,
  stringstyle=\color{attackRed},
  breaklines=true, frame=single, rulecolor=\color{black!30},
  showstringspaces=false, columns=fullflexible, tabsize=2,
  xleftmargin=10pt, framexleftmargin=8pt}

%---------------------------------------------------------------------
%  Macros matemáticos
%---------------------------------------------------------------------
\newcommand{\Hentro}{H}

%=====================================================================
\begin{document}
%=====================================================================

%---------------------------------------------------------------------
%  Portada
%---------------------------------------------------------------------
\begin{titlepage}
\centering
\vspace*{1.2cm}

\begin{tikzpicture}
  \node[rounded corners=2pt, fill=itbaCyan, text=black, inner sep=14pt,
        minimum width=0.8\textwidth, align=center] (t)
        {\Huge\bfseries Criptografía y Seguridad};
\end{tikzpicture}

\vspace{0.6cm}
{\LARGE Seguridad en aplicaciones}\\[0.25cm]
{\Huge\bfseries\color{itbaCyanDark} Flujo de Información y Malware}

\vspace{0.8cm}

% Ilustración: bóveda / caja fuerte (recreación del icono de portada)
\begin{tikzpicture}[scale=1.0]
  \fill[black!12] (-3.2,-2.4) rectangle (3.2,2.8);
  \draw[line width=2pt, black!55] (-3.2,-2.4) rectangle (3.2,2.8);
  \fill[itbaCyanDark!75] (-2.4,-2.0) rectangle (2.4,2.4);
  \draw[line width=3pt, black!70] (-2.4,-2.0) rectangle (2.4,2.4);
  \draw[line width=1.2pt, itbaCyan] (-2.0,-1.6) rectangle (2.0,2.0);
  \draw[line width=2pt, black!70] (0,-2.0) -- (0,2.4);
  \foreach \a in {0,45,...,315}{
    \draw[line width=2pt, black!70] (-1.1,0.2) -- ++(\a:0.6);
  }
  \draw[line width=2.5pt, black!70] (-1.1,0.2) circle (0.42);
  \fill[itbaCyan] (-1.1,0.2) circle (0.16);
  \draw[line width=2.5pt, black!70] (1.1,0.2) circle (0.42);
  \fill[itbaCyan] (1.1,0.2) circle (0.16);
  \foreach \y in {-1.2,-0.4,0.8,1.6}{
     \fill[black!70] (-2.2,\y) circle (0.07);
     \fill[black!70] (2.2,\y) circle (0.07);
  }
\end{tikzpicture}

\vfill
{\large Apunte integral — Teoría + Práctica}\\[0.2cm]
{\normalsize Clase 9 \quad·\quad Capítulos 16--17, \emph{Computer Security: Art and Science} (M.~Bishop)}\\[0.4cm]
{\normalsize Práctica: \emph{CWE Top 25} (MITRE / SANS) — Debilidades de software más peligrosas}\\[0.4cm]
{\small Instituto Tecnológico de Buenos Aires (ITBA)}

\vspace{0.6cm}
{\footnotesize\itshape Documento elaborado a partir de las transcripciones de la clase
teórica (Prof.\ Pablo Eduardo Abad) y de la clase práctica, junto con las diapositivas oficiales de ambas.}
\end{titlepage}

\tableofcontents
\newpage

%=====================================================================
\section{Introducción y contexto}
%=====================================================================

Esta clase cierra el bloque de \textbf{Seguridad en aplicaciones} con dos temas: en la
parte teórica, \textbf{flujo de información} (su última pata: confinamiento, canales
encubiertos y control de flujo) y \textbf{malware}; en la práctica, los \textbf{errores de
software más peligrosos} catalogados por el \emph{CWE Top 25} de MITRE/SANS.

\begin{slidebox}
\textbf{Hilo conductor de las clases previas.} ``Las clases anteriores estuvimos viendo
distintos temas de seguridad en aplicaciones desde la parte más teórica de políticas y
modelos, hasta problemas más concretos como autenticación y control de acceso. Vamos a
retomar un problema que vimos de manera informal al principio: el de \emph{flujo de
información}.''
\end{slidebox}

El \textbf{flujo de información} ya se introdujo en la Clase 8 (junto con los principios de
diseño seguro y las vulnerabilidades). Esta clase lo \emph{profundiza} con el tratamiento
formal —entropía de Shannon y seguimiento del flujo— y le agrega lo nuevo: \textbf{canales
encubiertos}, \textbf{ataques de canal lateral}, \textbf{técnicas de aislamiento} y, finalmente,
la familia completa del \textbf{malware}.

\begin{defbox}{Tríada CIA — el marco de referencia}
Todo lo que sigue (flujo de información, malware, errores de software) se evalúa por su
impacto sobre las tres propiedades clásicas:
\begin{itemize}[leftmargin=1.4em,itemsep=2pt]
  \item \textbf{Confidencialidad}: protege las \emph{lecturas}.
  \item \textbf{Integridad}: protege las \emph{modificaciones}.
  \item \textbf{Disponibilidad}: que quien deba acceder pueda hacerlo cuando lo necesita.
\end{itemize}
La práctica de esta clase usa explícitamente la tríada para clasificar el \emph{impacto} de
cada debilidad de software.
\end{defbox}

%=====================================================================
\section{El problema: control de acceso vs.\ flujo de información}
%=====================================================================

\begin{slidebox}
\textbf{El problema.} Una política de seguridad dice: ``\emph{Un empleado no puede ver los
sueldos de otros empleados}.'' \quad\textbf{¿Cómo podemos implementar esto?}
\end{slidebox}

El control de acceso parece la respuesta inmediata (es un problema de confidencialidad). Pero
el control de acceso \textbf{limita el acceso a operaciones sobre objetos}, no el destino de
la información una vez leída. Como los objetos \emph{contienen} información, limitar el acceso
limita la información\dots\ \emph{mientras la información sea estática}.

\begin{warnbox}{El control de acceso no impide copiar la información}
La información \textbf{no es estática}: es actualizada y \textbf{puede copiarse}. Un usuario
con permiso de \emph{leer} su sueldo podría escribirlo en otro lugar (un archivo de notas) que
otros usuarios sí pueden leer. El control de acceso, sentencia a sentencia, no detecta ese
\textbf{camino de lecturas y escrituras}.
\end{warnbox}

% ---- Diagrama: leer sueldo, escribir notas (recreación slide 3) ----
\begin{figure}[h]
\centering
\begin{tikzpicture}[font=\small,>=Stealth,node distance=8mm]
  \node[circle,draw,fill=itbaBlue,minimum size=0.9cm] (u) at (0,0) {};
  \node[font=\scriptsize] at (0,-0.85) {Empleado};
  \node[draw,rounded corners,fill=boxBlue,minimum width=2cm,minimum height=1cm] (sueldo) at (4,1) {SUELDO};
  \node[draw,rounded corners,fill=slideGray,minimum width=2cm,minimum height=1cm] (notas) at (4,-1) {NOTAS};
  \node[circle,draw,fill=itbaBlue,minimum size=0.7cm] (u2) at (6.2,-1) {};
  \draw[->,line width=1pt] (u) -- node[above,font=\scriptsize]{LEER} (sueldo);
  \draw[->,line width=1pt] (u) -- node[below,font=\scriptsize]{ESCRIBIR} (notas);
  \draw[->,line width=1pt,attackRed] (notas) -- node[above,font=\scriptsize]{LEER} (u2);
\end{tikzpicture}
\caption{La política se respeta sentencia a sentencia, pero la información \emph{fluye} de
SUELDO hacia un tercero vía NOTAS. \textit{(Recreación de la diapositiva 3.)}}
\end{figure}

\begin{defbox}{Políticas y control de acceso}
Las políticas, por lo general, \textbf{restringen el flujo de información}, no el acceso a
objetos. Comparar:
\begin{center}
\begin{tabular}{p{6.2cm} c p{6.2cm}}
\textbf{Evitar que un empleado \emph{sepa} el sueldo de otro} & vs. &
\textbf{Evitar que un empleado \emph{acceda} a la base de datos de sueldos}\\
\end{tabular}
\end{center}
\textbf{¿Entonces los controles de acceso no sirven?} Sí sirven, pero en términos formales son
\textbf{mecanismos abiertos}: aseguran una \emph{parte} del problema, no todo. \textbf{Deben
ser complementados.}
\end{defbox}

\begin{practbox}{La metáfora de las capas (aporte de la clase)}
Como los mecanismos son abiertos, las empresas piensan la seguridad como \textbf{capas de
cebolla} (o, como surgió en clase, el ``\emph{modelo del queso suizo}''): cada capa tiene
agujeros, pero al \textbf{superponer muchas} se espera que ningún agujero quede alineado de
punta a punta, y cada capa tape los huecos de la anterior.
\end{practbox}

%=====================================================================
\section{El problema del confinamiento}
%=====================================================================

\begin{defbox}{Confinamiento}
El \textbf{problema del confinamiento (de acceso)} se refiere a la dificultad de asegurar que
un programa o proceso (\emph{sujeto}) \textbf{no transmita información} a otra entidad
(\emph{objeto}) a la que no está autorizado a acceder, ya sea de manera \textbf{directa} o
\textbf{indirecta}.
\end{defbox}

Si se pudiera \textbf{confinar el flujo de información dentro de un ambiente controlado}, no
habría riesgos de confidencialidad ni de integridad. El modelo \textbf{Bell--LaPadula} (visto
en la Clase 6) define justamente un modelo de seguridad para asegurar el confinamiento de la
confidencialidad: en un ambiente totalmente cerrado donde se aplica el 100\,\% del modelo, da
garantías \emph{respecto del flujo}, no solo del lugar donde está la información.

% ---- Reticulado Bell-LaPadula (recreación slide 6) ----
\begin{figure}[h]
\centering
\begin{tikzpicture}[font=\scriptsize,>=Stealth,node distance=6mm,
   lbl/.style={draw,rounded corners,fill=itbaBlue!25,inner sep=3pt}]
  \node[lbl] (top) at (0,3) {(Top Secret, \{nuclear, crypto\})};
  \node[lbl] (tn) at (-4.6,1.5) {(Top Secret, \{nuclear\})};
  \node[lbl] (snc) at (0,1.5) {(Secret, \{nuclear, crypto\})};
  \node[lbl] (tc) at (4.6,1.5) {(Top Secret, \{crypto\})};
  \node[lbl] (sn) at (-4.6,0) {(Secret, \{nuclear\})};
  \node[lbl] (t0) at (0,0) {(Top Secret, \{\})};
  \node[lbl] (sc) at (4.6,0) {(Secret, \{crypto\})};
  \node[lbl] (s0) at (0,-1.5) {(Secret, \{\})};
  \draw (top)--(tn); \draw (top)--(snc); \draw (top)--(tc);
  \draw (tn)--(sn); \draw[dashed] (tn)--(t0);
  \draw (snc)--(sn); \draw[dashed] (snc)--(t0); \draw (snc)--(sc);
  \draw[dashed] (tc)--(t0); \draw (tc)--(sc);
  \draw (sn)--(s0); \draw (t0)--(s0); \draw (sc)--(s0);
\end{tikzpicture}
\caption{Reticulado de etiquetas de seguridad (nivel, categorías) de Bell--LaPadula que ordena
el flujo permitido. \textit{(Recreación de la diapositiva 6.)}}
\end{figure}

\subsection{Modelo ideal: aislación total}
El modelo ideal de confinamiento es la \textbf{aislación total}. Sus requisitos son simples:
\begin{enumerate}[leftmargin=1.6em,itemsep=2pt]
  \item Un proceso \textbf{no puede comunicarse} con ningún otro proceso.
  \item Un proceso \textbf{no puede ser observado} por nadie.
\end{enumerate}
\textbf{Consecuencia}: el proceso no revela información. ``El sistema más seguro del mundo corre
en una computadora apagada, metida en una caja fuerte, enterrada en el fondo del océano, sin que
nadie sepa dónde está.''

\begin{warnbox}{El problema de la aislación total}
La aislación total \textbf{le quita utilidad al sistema}: no le podemos dar entrada ni observar
la salida. Y, en la práctica, lo más difícil es el \emph{segundo} requisito: \textbf{cualquier
cosa que corra en una computadora tiene medidas observables}, porque usa recursos:
\begin{center}
\begin{tabular}{ll}
\textbf{Memoria} & \textbf{Ciclos de CPU}\\
\textbf{Espacio en disco} & \textbf{Ancho de banda}\\
\end{tabular}
\end{center}
Observando esos recursos se puede \textbf{extraer información} del proceso y de su resultado.
\end{warnbox}

%=====================================================================
\section{Canales encubiertos (Covert Channels)}
%=====================================================================

\begin{defbox}{Canal encubierto (covert channel)}
Un \textbf{canal encubierto} es un método de comunicación usado para transferir información de
manera \textbf{oculta}, aprovechando \textbf{vías no diseñadas para transmitir datos}. En el
contexto de seguridad son un problema porque, al no haber sido pensados para transmitir, es muy
poco probable que tengan los \textbf{controles} que sí tienen los canales legítimos: sirven para
\textbf{evadir los controles de seguridad y transmitir datos sin detección}.
\end{defbox}

\subsection{Canales de temporización (Timing Channels)}
Codifican datos en la \textbf{variación del tiempo entre eventos}. Se monitorea el canal a
intervalos regulares y, según cómo cambie la lectura, se infiere información. Ejemplos:
\begin{itemize}[leftmargin=1.4em,itemsep=1pt]
  \item \textbf{Latencia de red}, \textbf{carga de CPU}, \textbf{tiempos de respuesta de
        aplicaciones}.
\end{itemize}

\begin{practbox}{Traffic shaping (aporte de la clase)}
Sin entender el \emph{contenido}, contando la cantidad de tramas que circulan se puede inferir
\emph{cuándo} hay comunicación y hasta \emph{qué tipo} (un \emph{request/response}, un
\emph{streaming} multimedia). Los proveedores de Internet lo usan para \textbf{traffic
shaping}: analizan la \emph{forma} de la trama y aplican \emph{throttling} (p.\ ej.\ a peer-to-peer
o video) cuando se excede el ancho de banda. Son aplicaciones ``grises'': no maliciosas, pero la
misma técnica permite \emph{extraer información}.
\end{practbox}

\subsection{Canales de almacenamiento (Storage Channels)}
Usan el \textbf{almacenamiento de información en áreas no convencionales o no reguladas} del
sistema:
\begin{itemize}[leftmargin=1.4em,itemsep=1pt]
  \item \textbf{Espacio en archivo}: bits no usados / espacio de relleno (\emph{padding}).
  \item \textbf{Atributos de archivos}: metadata como fecha de creación, modificación o acceso.
  \item \textbf{Bits no usados en estructuras de datos}: bits reservados, o lugares de protocolo
        donde ``debe ir un valor aleatorio'' y el valor tal vez no sea tan aleatorio.
\end{itemize}

\subsection{Ejemplos de canales encubiertos}
\begin{slidebox}
\begin{itemize}[leftmargin=1.4em,itemsep=2pt]
  \item \textbf{Exfiltración de información codificando los sonidos del proceso de impresión}
        (un micrófono de alta fidelidad reconstruye lo que se imprime).
        \href{https://dl.acm.org/doi/10.1145/3510583}{\texttt{dl.acm.org/doi/10.1145/3510583}}
  \item \textbf{ICMP \& DNS Tunneling}.
        \href{https://www.akamai.com/blog/news/introduction-to-dns-data-exfiltration}{Akamai —
        introduction to DNS data exfiltration}
  \item \textbf{Emanaciones electromagnéticas}.
\end{itemize}
\end{slidebox}

\begin{practbox}{DNS tunneling explicado (aporte de la clase)}
Muchas empresas filtran el tráfico saliente con \emph{firewalls}, pero \textbf{casi todo
Internet necesita DNS} para resolver nombres, y durante años los firewalls dejaban pasar las
consultas DNS sin control. La idea: si controlás un dominio y su servidor DNS, desde la máquina
víctima hacés búsquedas de \textbf{subdominios} cuyo nombre es —p.\ ej.— un bloque del dato a
exfiltrar codificado en Base64. Cada consulta llega a tu servidor DNS y, juntando los nombres de
los subdominios, \textbf{se rearma información arbitraria}, evadiendo los controles. Los canales
encubiertos creativos son ``bastante complicados de evitar''.
\end{practbox}

%=====================================================================
\section{Ataques de canal lateral (Side-Channel Attacks)}
%=====================================================================

\begin{defbox}{Ataque de canal lateral}
Un \textbf{ataque de canal lateral} explota las \textbf{emisiones físicas y temporales} de un
sistema para extraer información sensible. \textbf{No} se basa en vulnerabilidades clásicas del
software ni en errores de autenticación, sino en la \textbf{observación de efectos colaterales
del hardware} durante su funcionamiento. Es especialmente relevante en criptografía: midiendo
\textbf{tiempos de ejecución, consumo de energía, radiación electromagnética o ruido acústico}
se pueden inferir \textbf{claves de cifrado}.
\end{defbox}

\begin{warnbox}{Por qué es devastador en criptografía}
Toda la seguridad montada sobre criptografía depende de un eslabón: la \textbf{clave}. Si un
ataque de canal lateral recupera la clave, ``se cae toda la seguridad como un castillo de
naipes''. Aun \textbf{sin éxito total}, basta \emph{sesgar} la probabilidad de algunos bits:
una clave de 128 bits con 40 bits sesgados deja al atacante con $2^{88}$ pruebas por fuerza
bruta en lugar de $2^{128}$.
\end{warnbox}

\subsection{Ejemplos}
\begin{slidebox}
\begin{itemize}[leftmargin=1.4em,itemsep=2pt]
  \item \textbf{Timing attack}: medir el tiempo de operaciones criptográficas para inferir la
        clave. (La cantidad de ciclos de reloj para multiplicar/exponenciar números grandes
        \emph{depende del tamaño} de esos números; combinado con texto plano conocido, se va
        despejando la incógnita.)
  \item \textbf{Power analysis}: analizar el \textbf{consumo de energía} del dispositivo durante
        las operaciones criptográficas.
  \item \textbf{Emanaciones electromagnéticas}: usar la radiación EM para inferir la operación
        interna y las claves.
\end{itemize}
\end{slidebox}

\begin{practbox}{La familia moderna: Rowhammer, Spectre, Meltdown (aporte de la clase)}
\begin{itemize}[leftmargin=1.4em,itemsep=1pt]
  \item \textbf{Rowhammer}: golpea (lee/escribe) repetidamente filas de memoria para alterar
        \emph{por física} bits vecinos sin acceso directo; uno de los más difíciles de evitar.
  \item \textbf{Spectre / Meltdown}: explotan la ejecución especulativa de los procesadores.
\end{itemize}
Ganaron mucha importancia con los \textbf{hyperscalers} (AWS, Google, Azure): en una misma
máquina física conviven servidores virtuales de \emph{distintos clientes}. Si desde una VM se
pudiera leer la memoria de otra VM en el mismo host, ``es una lotería, pero se corre sin parar
hasta que sale algo interesante''. Las \textbf{smart cards} (transporte, acceso a edificios) son
otro blanco clásico: tienen memoria/claves cifradas, pero como están en manos del atacante,
históricamente se las rompió midiendo \textbf{diferencias de consumo}.
\end{practbox}

% ---- Comparación covert vs side-channel (recreación slide 16) ----
\begin{figure}[h]
\centering
\begin{tikzpicture}[font=\scriptsize,>=Stealth]
  % Canal encubierto
  \node[draw,rounded corners,fill=practGreenL,minimum width=2cm,align=center] (unt) at (-1.5,0.8) {Proceso\\ no confiable};
  \node[draw,rounded corners,fill=itbaBlue!30,minimum width=2cm,align=center] (sens) at (2.5,0.8) {Información\\ sensible};
  \node[circle,draw=attackRed,line width=1.5pt,fill=white,minimum size=0.7cm] (block) at (0.5,0.8) {\textcolor{attackRed}{\Large$\oslash$}};
  \node[font=\scriptsize] at (0.5,1.55) {Security Policy};
  \node[cloud,draw,fill=slideGray,cloud puffs=11,minimum width=2.2cm,minimum height=1.1cm,aspect=1.6] (cc) at (0.5,-0.9) {Covert Channel};
  \draw[<->,gray] (unt) -- (block); \draw[<->,gray] (block) -- (sens);
  \draw[->,gray] (unt) to[bend right=20] (cc);
  \draw[->,gray] (cc) to[bend right=20] (sens);
  \node[align=center,font=\scriptsize] at (0.5,-2.3) {\textbf{Canal encubierto:}\\ evade controles para transmitir\\ información oculta};
  % Side channel
  \node[draw,rounded corners,fill=itbaBlue!30,minimum size=0.7cm] (chip) at (7,0.8) {chip};
  \node[draw,rounded corners,minimum width=1.2cm] (probe) at (8.8,0.8) {$\sim$ medición};
  \node[draw,rounded corners,fill=attackRed!12,minimum width=1.2cm] (att) at (8.8,-0.6) {atacante};
  \node[circle,draw=attackRed,minimum size=0.6cm,fill=attackRed!15] (key) at (7,-0.6) {\textcolor{attackRed}{clave}};
  \draw[->,attackRed,line width=1pt] (chip) -- node[above,font=\tiny]{emisión} (probe);
  \draw[->,gray] (probe) -- (att);
  \draw[->,attackRed,line width=1pt] (att) -- (key);
  \node[align=center,font=\scriptsize] at (7.9,-2.3) {\textbf{Canal lateral:}\\ observa el comportamiento físico\\ del hardware en operación};
\end{tikzpicture}
\caption{\textbf{Canal encubierto} (evadir controles para transmitir) vs.\ \textbf{ataque de
canal lateral} (recolectar información por la observación física del hardware).
\textit{(Recreación de la diapositiva 16.)}}
\end{figure}

%=====================================================================
\section{Control de flujo}
%=====================================================================

Hay dos grandes familias de técnicas para hacer cumplir (\emph{enforce}) el control de flujo:

\begin{defbox}{Compile-time vs.\ run-time enforcement}
\begin{description}[leftmargin=2.2em,style=nextline,itemsep=3pt]
  \item[Tiempo de compilación] Analiza la \textbf{propagación de información directamente en el
        código fuente}. Es \textbf{costoso y complejo} $\Rightarrow$ limitado a componentes muy
        sensibles.
  \item[Tiempo de ejecución] \textbf{Control dinámico} (sandboxing) y \textbf{aislación}
        (virtualización). \emph{Son los métodos más utilizados.}
\end{description}
\end{defbox}

%---------------------------------------------------------------------
\subsection{Control de flujo en tiempo de compilación}
%---------------------------------------------------------------------

\textbf{Concepto}: cuantificar la información asociada a cada variable, evaluar cómo cambia en
cada sentencia, y comparar el valor inicial y final. Si hay una \textbf{transferencia que no
debiera ocurrir}, fallar.

\subsubsection{Entropía como medida de información}
Hace casi 90 años, \textbf{Claude Shannon} propuso convertir en un solo número la cantidad de
información que puede circular por un canal discreto.

\begin{defbox}{Entropía de Shannon}
Sea $X$ una variable aleatoria discreta (V.A.D.) que toma valores $x_1,\dots,x_n$:
\[
  \Hentro(X) \;=\; -\sum_{i=1}^{n} p(x_i)\,\lg p(x_i)
\]
(el $\lg$ es el logaritmo en base 2; el resultado se mide en \textbf{bits}).
\textbf{Mide la incertidumbre} a la hora de determinar el valor de la variable: la cantidad de
información que \emph{nos falta} conocer para describir completamente el canal.
\begin{itemize}[leftmargin=1.4em,itemsep=1pt]
  \item \textbf{Máximo}: $p(x_i)=1/n$ (distribución \textbf{uniforme}) — no sabemos nada, puede
        venir cualquier valor.
  \item \textbf{Mínimo} ($=0$): $p(x_i)=1,\ p(x_j)=0$ para $j\neq i$ (\textbf{evento único}) —
        sabemos exactamente qué se transmite, no hay información extra.
\end{itemize}
\end{defbox}

Es una función \textbf{continua}: cuantifica también los escenarios intermedios, donde
conocemos \emph{algo} de la distribución pero no todo.

\begin{defbox}{Entropía condicional}
Sean $X$ ($x_1,\dots,x_n$) e $Y$ ($y_1,\dots,y_m$) dos V.A.D. Para un valor concreto de $Y$:
\[
  \Hentro(X \mid Y=y) \;=\; -\sum_{i=1}^{n} p(x_i \mid y)\,\lg p(x_i \mid y)
\]
y, en general, como promedio ponderado de las entropías condicionales:
\[
  \Hentro(X \mid Y) \;=\; \sum_{j=1}^{m} p(y_j)\, \Hentro(X \mid Y=y_j).
\]
\end{defbox}

\subsubsection{Definición formal de flujo de información}
\begin{defbox}{Flujo de información}
Sea $s$ el estado de un sistema y $t$ el estado luego de ejecutar comandos $c_1,\dots,c_n$.
Sean $x_s, y_s$ valores de objetos en el estado $s$, e $y_t$ el valor de $y$ en el estado $t$.
\textbf{Hay flujo de información de $x$ a $y$} si:
\[
  \Hentro(x_s \mid y_t) < \Hentro(x_s \mid y_s) \quad\text{si } y \text{ existe en } s,
\]
\[
  \Hentro(x_s \mid y_t) < \Hentro(x_s) \quad\text{si } y \text{ \emph{no} existe en } s.
\]
\end{defbox}

Es decir: si lo que \emph{desconozco} de $x$ \textbf{disminuye} al conocer el estado final de
$y$, entonces hubo transferencia de información de $x$ a $y$. Si no hubiera transferencia,
debería dar \textbf{mayor o igual} ($\geq$): el ``mayor'' aparece, por ejemplo, cuando una
asignación \emph{pisa} la variable y destruye información.

\begin{slidebox}
\textbf{En todo programa hay flujo de información} —es necesario para funcionar—. Por eso el
seguimiento se \textbf{ciñe a casos puntuales críticos}, típicamente: verificar que una
\textbf{función de cifrado no revele información de la clave}.
\end{slidebox}

\subsubsection{Flujo directo (ejercicio de la slide)}
\begin{practbox}{Ejercicio resuelto — flujo directo}
Considerar el comando \texttt{y = x + z}, donde $0 \le x \le 7$ con igual probabilidad
($8$ valores) y $Z=\{p(z{=}1)=0.5,\ p(z{=}2)=0.25,\ p(z{=}3)=0.25\}$.

\medskip
\textbf{Antes del comando} (entropía de $x$, distribución uniforme sobre $8$ valores):
\[
  \Hentro(x) = -\sum_{i=1}^{8} \tfrac{1}{8}\lg\tfrac{1}{8}
             = -8\cdot\tfrac{1}{8}\lg\tfrac{1}{8}
             = \lg 8 = \boxed{3}.
\]
\textbf{Conociendo $y$ luego del comando}: $x$ solo puede tomar $3$ valores ($y{-}1,\,y{-}2,\,y{-}3$),
con las probabilidades heredadas de $z$:
\[
  \Hentro(x \mid y) = -\tfrac{1}{2}\lg\tfrac{1}{2} - 2\cdot\tfrac{1}{4}\lg\tfrac{1}{4}
                    = \tfrac{1}{2} + 2\cdot\tfrac{1}{2} = \boxed{1.5}.
\]
Como $1.5 < 3$, \textbf{hay traspaso de información de $x$ a $y$}.
\end{practbox}

\subsubsection{Flujo indirecto}
La información puede fluir sin que las variables aparezcan en la misma asignación.

\begin{practbox}{Ejercicio resuelto — flujo indirecto (condicional)}
\texttt{if (x == 0) y = 1; else y = 0;} \quad ($x,y$ \emph{nunca} aparecen en la misma asignación).
Con $x\in\{0,1\}$ y $p(x{=}0)=0.5$:
\[
  \Hentro(x) = -2\cdot\tfrac{1}{2}\lg\tfrac{1}{2} = 1, \qquad
  \Hentro(x \mid y) = 0.
\]
Como $0 < 1$, \textbf{hay traspaso de información} (observando $y$ se deduce $x$). Es un
\textbf{canal espacial}.
\end{practbox}

\begin{practbox}{Ejercicio resuelto — flujo indirecto por comportamiento (loop)}
\texttt{while (x == 0) \{\}} \quad (no existe ninguna asignación). Con $x\in\{0,1\}$,
$p(x{=}0)=0.5$. Definimos una variable observable $y=0$ \emph{si el programa termina}:
\[
  \Hentro(x) = 1, \qquad \Hentro(x \mid y) = 0 \quad\Rightarrow\quad \textbf{hay traspaso}.
\]
Observar \emph{si el programa termina o no} (un \textbf{canal temporal por comportamiento}) da
información indirecta sobre $x$.
\end{practbox}

\subsubsection{Límites}
Estas técnicas comparten los límites de la \textbf{verificación formal} de sistemas:
\begin{itemize}[leftmargin=1.4em,itemsep=1pt]
  \item \textbf{No es computable verificar todo programa} (límite teórico; en general los no
        computables son los que nunca terminan).
  \item Cuando hay \textbf{interactividad} (input externo), los casos \textbf{explotan} en
        cantidad de forma exponencial.
  \item \textbf{Los ciclos son un problema}: se hace \emph{loop unrolling} (desenrollar el ciclo
        $N$ veces); analizar hasta $\sim 10$ iteraciones suele bastar porque los problemas
        estructurales aparecen en pocos ciclos.
\end{itemize}

%---------------------------------------------------------------------
\subsection{Control de flujo en tiempo de ejecución: aislamiento}
%---------------------------------------------------------------------

\begin{defbox}{Las dos formas de aislar procesos}
\begin{enumerate}[leftmargin=1.6em,itemsep=2pt]
  \item \textbf{Virtualización}: simularle a la aplicación un \textbf{ambiente de ejecución
        sintético}, separado del real y aislado por definición.
  \item \textbf{Sandboxing}: \textbf{controlar y modificar las interacciones} del sistema con su
        entorno para aislarlo.
\end{enumerate}
\end{defbox}

\subsubsection{Virtualización (Virtual Machines)}
Las VMs usan \textbf{hipervisores} para emular hardware y gestionar distintos sistemas
operativos sobre un solo host físico:
\begin{itemize}[leftmargin=1.4em,itemsep=1pt]
  \item \textbf{Bare-metal} (tipo 1): KVM, VMware ESXi, Hyper-V.
  \item \textbf{Hosted} (tipo 2): VirtualBox, VMware Workstation.
\end{itemize}
Cada VM tiene su \textbf{propio kernel} $\Rightarrow$ aislamiento fuerte (ni siquiera comparten
CPU lógicamente). Contrapartida: es \textbf{pesadísimo}, hay que emular todo el ambiente.

\subsubsection{Sandboxing}
Analiza las interacciones de un proceso con su ambiente y \textbf{las modifica para aislarlo}.
Ejemplo de diseño: la \textbf{isolación de sitios} de Chromium, donde cada \emph{renderer}
(\texttt{a.com}, \texttt{b.com}, \texttt{c.com}) corre en su propio sandbox sobre un único
\emph{browser process}.

\begin{practbox}{\texttt{chroot} y arbitraje de E/S (aporte de la clase)}
En UNIX \textbf{toda entrada/salida es un archivo}. El comando \texttt{chroot} lanza un proceso
con una \textbf{vista cortada del filesystem}: la carpeta que le indiquemos pasa a ser su raíz
\texttt{/}. Uso clásico: los \textbf{instaladores de paquetes} corren los scripts de terceros
dentro de un sandbox (una vista limpia con sus \texttt{/etc}, \texttt{/bin}, etc.); luego se
verifica \emph{qué archivos nuevos aparecieron} y, si nada es sospechoso, recién ahí se copian
al sistema real. A diferencia de las VMs, en sandboxing los procesos \textbf{comparten memoria y
CPU} $\Rightarrow$ ahí quedan canales no aislados.
\end{practbox}

\paragraph{Aislamiento a nivel sistema operativo.}
\begin{itemize}[leftmargin=1.4em,itemsep=1pt]
  \item \textbf{Linux}: \emph{namespaces} (aíslan procesos, usuarios, red, filesystem\dots) y
        \emph{cgroups} (controlan/limitan CPU, memoria y red por grupo de procesos).
  \item \textbf{Windows}: \emph{User Mode Scheduling} (similar a namespaces) y \emph{Job Objects}
        (gestionan grupos de procesos como una unidad, con límites de recursos y prioridades).
\end{itemize}

\subsubsection{Contenedores}
Los \textbf{contenedores} aprovechan \emph{namespaces} y \emph{cgroups} (Docker es el ejemplo
popular) para aislar aplicaciones: múltiples contenedores comparten el \textbf{mismo kernel del
host} pero mantienen separación lógica de red, procesos y filesystem. Windows soporta
contenedores nativos usando características de Windows o Hyper-V. Eliminan la necesidad de un
SO completo por instancia (a diferencia de las VMs), dando casi la misma abstracción pero más
liviana.

% ---- VM vs Container vs Bare metal (recreación slide 40) ----
\begin{figure}[h]
\centering
\begin{tikzpicture}[font=\scriptsize,
  layer/.style={draw,minimum width=3cm,minimum height=0.5cm},
  app/.style={draw,minimum width=0.95cm,minimum height=0.5cm}]
  % Virtual server
  \node[font=\bfseries] at (0,3.5) {Virtual Server};
  \foreach \x/\c in {-1/itbaBlue,0/itbaBlue,1/itbaBlue}{
     \node[app,fill=\c!40] at (\x,3.0) {APP}; \node[app,fill=\c!25] at (\x,2.5) {bin/lib};
     \node[app,fill=\c!55] at (\x,2.0) {OS};}
  \node[layer,fill=attackRed!30] at (0,1.5) {Hypervisor};
  \node[layer,fill=cweProtect] at (0,1.0) {Host Operating System};
  \node[layer,fill=black!10] at (0,0.5) {Hardware};
  % Container
  \node[font=\bfseries] at (5,3.5) {Container};
  \foreach \x/\c in {4/practGreen,5/practGreen,6/practGreen}{
     \node[app,fill=\c!40] at (\x,3.0) {APP}; \node[app,fill=\c!25] at (\x,2.5) {bin/lib};}
  \node[layer,fill=attackRed!30] at (5,2.0) {Container Engine};
  \node[layer,fill=cweProtect] at (5,1.5) {Host Operating System};
  \node[layer,fill=black!10] at (5,1.0) {Hardware};
  % Bare metal
  \node[font=\bfseries] at (10,3.5) {Bare Metal};
  \foreach \x/\c in {9/cweOther,10/cweOther,11/cweOther}{
     \node[app,fill=\c] at (\x,3.0) {APP}; \node[app,fill=black!8] at (\x,2.5) {bin/lib};}
  \node[layer,fill=cweProtect] at (10,2.0) {Host Operating System};
  \node[layer,fill=black!10] at (10,1.5) {Hardware};
\end{tikzpicture}
\caption{\emph{Stack} comparativo: máquina virtual (cada app con su propio OS sobre el
hipervisor), contenedor (apps comparten el kernel del host vía el \emph{container engine}) y
\emph{bare metal}. \textit{(Recreación de la diapositiva 40.)}}
\end{figure}

%=====================================================================
\section{Malware}
%=====================================================================

\begin{defbox}{Malware (Malicious Software)}
\textbf{Malware} es el término con el que se identifica a cualquier programa \textbf{diseñado
para causar daño, atravesar un control o extraer información}. Un ataque \textbf{usualmente usa
múltiples tipos de malware} para ejecutar cada parte del ciclo del ataque, y suele
\textbf{aprovechar vulnerabilidades del software} para lograr su cometido.
\end{defbox}

\begin{slidebox}
El malware se ubica después de las vulnerabilidades en la cadena: hay \emph{amenazas} que se
materializan en \emph{vulnerabilidades}, pero hay otra familia —\textbf{programas diseñados
específicamente con un fin malicioso}—. Lo que les da el nombre distinto no es la técnica
(siguen siendo programas), sino que \textbf{fueron concebidos exclusivamente con un fin
malicioso}.
\end{slidebox}

% ---- Taxonomía de malware (diagrama propio) ----
\begin{figure}[h]
\centering
\begin{tikzpicture}[font=\small,>=Stealth,
  cat/.style={draw,rounded corners,fill=attackRed!12,minimum width=2.5cm,minimum height=0.7cm,align=center},
  root/.style={draw,rounded corners,fill=itbaCyan,minimum width=2.2cm,minimum height=0.8cm}]
  \node[root] (m) at (0,0) {\bfseries Malware};
  \node[cat] (vw) at (-5.2,-1.8) {Virus / Worm\\ \scriptsize (propagación)};
  \node[cat] (tr) at (-2.6,-1.8) {Troyano\\ \scriptsize (ocultamiento)};
  \node[cat] (rk) at (0,-1.8) {Rootkit\\ \scriptsize (persistencia oculta)};
  \node[cat] (sp) at (2.6,-1.8) {Spyware\\ \scriptsize (recolección)};
  \node[cat] (rw) at (5.2,-1.8) {Ransomware\\ \scriptsize (extorsión)};
  \node[cat] (bn) at (0,-3.3) {Botnet \scriptsize (zombis multiuso) — Command \& Control};
  \foreach \n in {vw,tr,rk,sp,rw} \draw[->] (m) -- (\n);
  \draw[->] (m.south) to[bend right=10] (bn.west);
\end{tikzpicture}
\caption{Taxonomía de malware. \emph{Virus/worm} aportan la propagación; \emph{troyano} y
\emph{rootkit}, las formas de esconderse; \emph{spyware}/\emph{ransomware}/\emph{botnet}
describen \emph{qué hace} (el \emph{payload}). Las categorías \textbf{se combinan}.}
\end{figure}

\begin{practbox}{Propagación vs.\ payload (aclaración del profe)}
\textbf{Virus} y \textbf{worm} describen una \emph{forma de propagarse}; \textbf{troyano} y
\textbf{rootkit}, una \emph{forma de esconderse}. Ninguno de los cuatro dice \emph{qué hace} una
vez adentro: eso se llama \textbf{payload}. Las cosas se combinan: un \emph{virus} se propaga
entre máquinas, en cada una se \emph{enquista} como \emph{rootkit}, y su \emph{payload} puede ser
un \emph{ransomware}.
\end{practbox}

\subsection{Virus / Worm}
\begin{itemize}[leftmargin=1.4em,itemsep=1pt]
  \item Característica clave: \textbf{infectar una máquina y propagarse} usando los medios
        disponibles. Puede infectar por \textbf{explotación de una vulnerabilidad remota} o
        reescribiendo un \textbf{recurso compartido} (archivo ejecutable).
  \item Un \textbf{virus} suele tener dos componentes: \textbf{replicación} + \textbf{ejecución}
        (según el componente de ejecución, se le da otro nombre).
  \item El \textbf{worm} se especializa en la \textbf{infección por red}, propagándose incluso
        por email (\textbf{Melissa}, 1999). El primer gusano nació por error en un experimento de
        cómputo distribuido y saturó las máquinas por falta de memoria.
  \item \textbf{SQL Slammer} infectó \textbf{75.000 computadoras en 10 minutos} aprovechando bases
        de datos SQL expuestas a Internet (asociados a \emph{denegación de servicio}).
\end{itemize}

\subsection{Troyano}
\begin{itemize}[leftmargin=1.4em,itemsep=1pt]
  \item \textbf{Malware oculto dentro de otro software} que parece benigno y de interés para la
        víctima (del caballo de Troya). Se usa para \textbf{establecer al atacante dentro del
        equipo infectado}.
  \item Típico en \emph{cracks} / versiones piratas (Windows, Office) y en videojuegos móviles
        que empresas compran y modifican para incluir troyanos.
  \item \textbf{PC-Write} (1986): uno de los primeros \emph{shareware} (editor de texto) que
        contenía un troyano que \textbf{borraba todos los archivos} de la máquina.
\end{itemize}

\subsection{Rootkit}
\begin{itemize}[leftmargin=1.4em,itemsep=1pt]
  \item Software diseñado para \textbf{ocultar malware} dentro de la máquina infectada, con
        \textbf{contramedidas activas}: ofuscación, encripción, reescritura y reemplazo de
        herramientas del sistema para quedar oculto a administradores y antivirus.
  \item Pueden operar a \textbf{nivel kernel o de usuario}. Un rootkit en Linux puede
        \textbf{interceptar system calls} (p.\ ej.\ ocultar su proceso de \texttt{ps} o sus
        archivos de \texttt{/proc}). Los más nuevos se meten en el \textbf{firmware/BIOS} (de ahí
        nacen \emph{Secure Boot} y el boot firmado), persistiendo aunque se formatee el disco.
  \item \textbf{NTRootkit} (1999): primer rootkit para Windows NT.
        \textbf{Stuxnet} (2010): usó rootkit para recolectar información por \emph{meses} antes de
        lanzar el ataque que afectó las plantas nucleares de Irán (atribuido a servicios de
        inteligencia; las máquinas estaban aisladas de Internet).
\end{itemize}

\begin{warnbox}{Por qué los rootkits son tan peligrosos}
Trabajan ``un nivel atrás'': leés memoria y no hay nada, leés disco y no hay nada. Cuanto más
sofisticado, de más cosas se esconde $\Rightarrow$ los más avanzados \textbf{no son detectables
por antivirus}. Contrapartida: por trabajar a tan bajo nivel, una actualización del SO puede
romperlos.
\end{warnbox}

\subsection{Spyware}
\begin{itemize}[leftmargin=1.4em,itemsep=1pt]
  \item Software diseñado para \textbf{recolectar información} del usuario infectado, manteniendo
        el \textbf{perfil más bajo posible}.
  \item Se especializa en \textbf{ocultarse} y usualmente establece alguna forma de
        \textbf{covert channel (túnel)} para exfiltrar; en el mercado hogareño (sin equipo de
        seguridad) suele usar canales normales.
  \item Algunas empresas lo desarrollan \textbf{a pedido de gobiernos} para vigilancia digital o
        espionaje (caso mencionado: \emph{Palantir}).
\end{itemize}

\subsection{Ransomware}
\begin{itemize}[leftmargin=1.4em,itemsep=1pt]
  \item Software para \textbf{bloquear el acceso} a la información, pidiendo \textbf{dinero}
        (rescate) para restaurarla. Habilitado por las \textbf{criptomonedas} (mecanismo de cobro
        descentralizado y difícil de rastrear).
  \item Variantes: cifrar la información y pedir rescate por la clave; o \textbf{exfiltrarla} y
        pedir rescate para \emph{no publicarla}.
  \item \textbf{CryptoLocker} (2013): enviado vía mail, infectaba al abrir el correo; se estima
        $\sim$U\$D~3M robados ese año. \textbf{WannaCry} (2017): aprovechó una vulnerabilidad
        conocida de Windows para \textbf{distribuirse en la red local}; $>$200K máquinas, costo
        estimado U\$D~4B. Es \textbf{pesadilla en el ámbito empresarial}: un servidor o base de
        datos comprometido pone en riesgo activos muy importantes.
\end{itemize}

\subsection{Botnet}
\begin{itemize}[leftmargin=1.4em,itemsep=1pt]
  \item Software para desplegar \textbf{bots / agentes multiuso} que quedan a la
        \textbf{escucha de instrucciones} (\textbf{máquinas zombi}). Se conectan a un servidor de
        \textbf{Command \& Control (C\&C)} y preguntan ``qué tengo que hacer''.
  \item Usos: \textbf{ataques distribuidos (DDoS)}, minería de cripto, accesos anónimos,
        \textbf{distribuir spam}, robar credenciales bancarias, usarlas de \emph{pivot} para otro
        ataque (cuando se rastrea, se llega a la zombi y no al original).
  \item Las infraestructuras C\&C son \textbf{sistemas distribuidos descentralizados, sin punto
        único de falla}, difíciles de dar de baja. \textbf{Srizbi} (2008): la botnet más grande
        documentada; responsable de $\sim$la \textbf{mitad del spam mundial} ($\sim$60B mails/día).
\end{itemize}

\begin{practbox}{La ``economía gris'' (aporte de la clase)}
Servicios no estrictamente ilegales que requieren ``algo no muy santo'' por detrás: envío de
spam evadiendo filtros de Gmail (vía computadoras hogareñas infectadas), o ``cómputo distribuido
en alquiler'' que en realidad son botnets. Ejemplo moderno: servicios que ofrecen acceso a las
APIs de Anthropic/OpenAI a un cuarto del precio, presumiblemente usando \textbf{API keys
robadas} con estos mecanismos de propagación + spyware.
\end{practbox}

% ---- Botnet architecture (recreación slide 42) ----
\begin{figure}[h]
\centering
\begin{tikzpicture}[font=\scriptsize,>=Stealth,node distance=5mm,
  bot/.style={circle,draw,fill=itbaBlue!40,minimum size=0.55cm}]
  \node[draw,rounded corners,fill=attackRed!15,align=center] (bm) at (0,0) {Botmaster};
  \node[bot] (b1) at (3,0.9) {bot};
  \node[bot] (b2) at (4.2,0.2) {bot};
  \node[draw,rounded corners,fill=slideGray] (cs) at (3,-0.4) {client/server};
  \node[bot] (b3) at (4.4,-1.0) {bot};
  \node[bot] (b4) at (2,-1.1) {bot};
  \draw[->] (bm) -- (3,0.2);
  \draw (b1)--(b2); \draw (b2)--(cs); \draw (cs)--(b3); \draw (cs)--(b4);
  \node[draw,rounded corners,fill=itbaCyan!25,align=left] (act) at (7.3,0) {Spam\\ DDoS\\ Infect systems\\ Phishing};
  \draw[->] (4.6,-0.1) -- (act);
  \node[draw,rounded corners,fill=cweOther,align=center] (tg) at (10.4,0) {Targets\\ \scriptsize (IoT, PCs,\\ servidores, móviles)};
  \draw[->] (act) -- (tg);
  \node[font=\scriptsize\bfseries] at (3.5,-2.0) {Command \& Control (P2P o cliente/servidor)};
\end{tikzpicture}
\caption{Arquitectura de una botnet: el \emph{botmaster} comanda nodos C\&C (P2P o
cliente/servidor) que ejecutan actividades maliciosas sobre los \emph{targets}.
\textit{(Recreación de la diapositiva 42.)}}
\end{figure}

\subsection{Ciclo de ataque del malware}
El ciclo de un malware suele tener, esquemáticamente, estas fases. Lo interesante es que
\textbf{toda la parte intermedia} (desde el reconocimiento hasta el bucle de escalamiento) es
\textbf{genérica} respecto del payload final $\Rightarrow$ hoy se venden \textbf{malware kits}
que ``productizan'' esa parte (aunque ser genérico también los hace más detectables, y hubo kits
que eran a su vez troyanos contra quien los usaba).

% ---- Ciclo de ataque (recreación slides 43-52) ----
\begin{figure}[h]
\centering
\begin{tikzpicture}[font=\scriptsize,>=Stealth,
  ph/.style={draw,fill=attackRed!25,minimum width=2.0cm,minimum height=0.8cm,align=center}]
  \node[ph] (r) at (0,0) {Recon\\ inicial};
  \node[ph] (c) at (2.3,0) {Compromiso\\ inicial};
  \node[ph] (e) at (4.6,0) {Establecerse\\ \scriptsize (foothold)};
  \node[ph] (es) at (6.9,0) {Escalar\\ privilegios};
  \node[ph] (ir) at (9.2,0) {Recon\\ interno};
  \node[ph,fill=attackRed!40] (cm) at (11.5,0) {Completar\\ misión};
  \foreach \a/\b in {r/c,c/e,e/es,es/ir,ir/cm} \draw[->,line width=1pt] (\a)--(\b);
  % loop maintain presence + move laterally
  \node[ph,fill=attackRed!15] (mp) at (6.9,1.4) {Mantener\\ presencia};
  \node[ph,fill=attackRed!15] (ml) at (9.2,1.4) {Movimiento\\ lateral};
  \draw[->,gray] (es) -- (mp); \draw[->,gray] (mp) -- (ml); \draw[->,gray] (ml) -- (ir);
\end{tikzpicture}
\caption{Ciclo de ataque del malware. El lazo \emph{escalar $\to$ mantener presencia $\to$
movimiento lateral $\to$ recon interno} se repite varias veces. \textit{(Recreación de las
diapositivas 43--52.)}}
\end{figure}

\begin{center}
\renewcommand{\arraystretch}{1.3}
\begin{tabular}{>{\bfseries}p{3.4cm} p{9.2cm}}
\toprule
Fase & Qué ocurre \\
\midrule
Reconocimiento inicial & Se identifica a la víctima y se recopila información: responsables,
puntos débiles, intereses; se mapean redes sociales y se establecen relaciones falsas.\\
Compromiso inicial & El malware se ``\emph{weaponiza}'' (empaqueta para ser ejecutado por la
víctima); con un ataque dirigido se compromete una o más máquinas. Puede durar varios días.\\
Establecerse (\emph{foothold}) & Establece comunicación con su \textbf{C\&C}, asegura
persistencia y almacenamiento; baja y compone nuevas herramientas. Suele durar minutos.\\
Escalar privilegios & Eleva los privilegios obtenidos usando vulnerabilidades locales/remotas.
Suele durar minutos.\\
Recon interno & Analiza comunicaciones, mapea otras máquinas/usuarios de interés; del
descubrimiento baja y compone nuevos malwares.\\
Movimiento lateral & Explotación remota de vulnerabilidades para incrementar la base infectada.\\
Mantener la presencia & Asegura persistir en las nuevas máquinas infectadas.\\
Completar misión & Ejecuta el \textbf{payload}: extrae/acumula información y la envía en horarios
de menor monitoreo; en ransomware, infecta el máximo de máquinas antes de activar el bloqueo.\\
\bottomrule
\end{tabular}
\end{center}

%=====================================================================
\section{Práctica: CWE Top 25 — Errores de software más peligrosos}
%=====================================================================

\begin{practbox}{De qué va la práctica}
La parte práctica trabaja sobre el \textbf{CWE Top 25} (\emph{Common Weakness Enumeration}) de
\textbf{MITRE/SANS}: el ranking de las \textbf{25 debilidades de software más peligrosas}. Son
debilidades \textbf{fáciles de encontrar y explotar}, que pueden permitir al adversario
\textbf{controlar el sistema, robar datos o impedir el funcionamiento}.
\end{practbox}

\subsection{Metodología del ranking}
\begin{itemize}[leftmargin=1.4em,itemsep=1pt]
  \item El \textbf{puntaje} se calcula principalmente en función de la \textbf{frecuencia} de
        aparición y del \textbf{impacto/daño}. No es lo mismo algo muy frecuente pero de bajo
        impacto, que algo frecuente \emph{y} de alto impacto.
  \item Se publica \textbf{anualmente desde $\sim$2010}. Cada año varias organizaciones reportan
        los problemas observados; MITRE los cataloga y, procesando $\sim$un año de datos, produce
        el ranking. \textbf{Fluctúa}: hay errores que entran, salen y \textbf{vuelven a aparecer}.
  \item Por cada error, la página
        (\href{http://cwe.mitre.org/top25}{\texttt{cwe.mitre.org/top25}}) detalla: \textbf{fase
        del ciclo de vida} donde se origina (requisitos / arquitectura y diseño / implementación),
        \textbf{lenguajes} donde aparece más, \textbf{consecuencias} (impacto en la tríada CIA),
        \textbf{probabilidad de explotación}, ejemplos y \textbf{casos reales observados}.
\end{itemize}

\begin{slidebox}
\textbf{Recomendación de lectura.} Conviene leer los errores \textbf{agrupados por categoría},
no en el orden del ranking: errores de la misma categoría tienen \textbf{soluciones similares},
así que entender uno ayuda a entender los demás (``si entendés SQL injection, XSS es parecido'').
\end{slidebox}

\subsection{Evolución de la clasificación}
\begin{center}
\renewcommand{\arraystretch}{1.3}
\begin{tabular}{p{4.5cm} p{8.2cm}}
\toprule
\textbf{2011 — 3 grandes grupos} & \textbf{2025 — ``Pillars'' (más desglosado)} \\
\midrule
\begin{minipage}[t]{4.4cm}\raggedright
$\bullet$ Interacción insegura entre componentes\\
$\bullet$ Manejo riesgoso de recursos\\
$\bullet$ Defensas abiertas / porosas
\end{minipage}
&
\begin{minipage}[t]{8.1cm}\raggedright
\emph{Improper Neutralization}; \emph{Improper Control of a Resource Through its Lifetime};
\emph{Protection Mechanism Failure}; \emph{Improper Access Control}; \emph{Improper Adherence to
Coding Standards}; \emph{Incorrect Calculation}; \emph{Improper Interaction Between Multiple
Correctly Behaving Entities}; \emph{Insufficient Control Flow Management}; \emph{Incorrect
Comparison}; \emph{Improper Check or handling of exceptional conditions}.
\end{minipage}\\
\bottomrule
\end{tabular}
\end{center}
Los \emph{pillars} de 2025 terminan ``cayendo'' en los tres grupos originales (más el cuarto). Las
categorías \textbf{más frecuentes} son la 1ª, la 2ª y la 4ª:
\textbf{neutralización inadecuada}, \textbf{control inapropiado de recursos} y \textbf{control de
acceso inapropiado}. \emph{(Referencias:} \href{https://cwe.mitre.org/data/definitions/1000.html}{\texttt{cwe.mitre.org/data/definitions/1000.html}}\emph{.)}

% ---- Diagrama: pilares CWE (categorías clave) ----
\begin{figure}[h]
\centering
\begin{tikzpicture}[font=\scriptsize,>=Stealth,
  pil/.style={draw,rounded corners,minimum width=3.2cm,minimum height=0.9cm,align=center}]
  \node[pil,fill=cweNeutral] (n) at (0,0) {\textbf{Neutralización inadecuada}\\ \scriptsize validar la entrada};
  \node[pil,fill=cweResource] (r) at (3.6,0) {\textbf{Control inapropiado}\\ \textbf{de recursos}\\ \scriptsize crear/usar/destruir};
  \node[pil,fill=cweAccess] (a) at (7.2,0) {\textbf{Control de acceso}\\ \textbf{inapropiado}\\ \scriptsize authN / authZ};
  \node[pil,fill=cweProtect] (p) at (10.6,0) {\textbf{Falla en mecanismos}\\ \textbf{de protección}};
  \node[pil,fill=cweOther,minimum width=10.5cm] (o) at (5.3,-1.5) {\textbf{Otros}: cálculos incorrectos (Integer Overflow) · condiciones excepcionales (NULL Pointer) · flujo de control (Race Condition) · credenciales hardcoded · uso incorrecto de criptografía};
\end{tikzpicture}
\caption{Pilares/categorías del CWE Top 25. Las tres primeras concentran la mayoría de las
debilidades del ranking. \textit{(Recreación de las diapositivas prácticas 3--5.)}}
\end{figure}

\subsection{El CWE Top 25 — 2025, agrupado por categoría}

\subsubsection{Neutralización inadecuada}
\begin{practbox}{Idea y mitigaciones}
\textbf{Idea}: \emph{asegurarse de que el input/output tenga ciertas propiedades antes de usarse}.
No se valida correctamente la entrada y con ella se construye un comando (SQL, sistema operativo,
HTML\dots). \textbf{Mitigaciones comunes} (las mismas para toda la categoría): \textbf{formatear
la entrada}, \textbf{listas blancas / listas negras} (lo permitido / lo prohibido) y
\textbf{validación de la entrada}.
\end{practbox}
\begin{center}
\renewcommand{\arraystretch}{1.25}
\begin{tabular}{>{\bfseries}c p{8.8cm} c}
\toprule
\# & Debilidad & Rank \\
\midrule
1 & Cross-site Scripting (XSS) & 1 \\
2 & SQL Injection & 2 \\
3 & OS Command Injection & 9 \\
4 & Code Injection & 10 \\
5 & Improper Input Validation & 18 \\
6 & Command Injection & 23 \\
\bottomrule
\end{tabular}
\end{center}

\subsubsection{Control inapropiado de un recurso}
\begin{practbox}{Idea}
\textbf{La creación, uso y destrucción de un recurso debe hacerse correctamente.} Muchos tienen
que ver con \textbf{buffer overflow} y manejo de memoria (más ligado a lenguajes de bajo nivel
como C); este año aparecen muchísimos. Se distingue el \emph{classic buffer overflow} (no se
chequea el tamaño de la entrada) del que se produce, p.\ ej., por llamadas recursivas sucesivas
(stack vs.\ heap). \emph{Path traversal}: armar un \emph{path} que termina dando acceso a cosas
que no se deberían tener. \emph{Unrestricted upload}: no controlar un archivo subido (un PNG que
en realidad es un ejecutable). \emph{Deserialization}: convertir una cadena/estructura sin
control.
\end{practbox}
\begin{center}
\renewcommand{\arraystretch}{1.25}
\begin{tabular}{>{\bfseries}c p{8.8cm} c}
\toprule
\# & Debilidad & Rank \\
\midrule
1 & Out-of-bounds Write (escritura fuera de límites) & 5 \\
2 & Path Traversal & 6 \\
3 & Use After Free & 7 \\
4 & Out-of-bounds Read (lectura fuera de límites) & 8 \\
5 & Buffer Copy without Checking Size of Input (Classic Buffer Overflow) & 11 \\
6 & Unrestricted Upload of File with Dangerous Type & 12 \\
7 & Stack-based Buffer Overflow & 14 \\
8 & Deserialization of Untrusted Data & 15 \\
9 & Heap-based Buffer Overflow & 16 \\
10 & Exposure of Sensitive Information to an Unauthorized Actor & 20 \\
11 & Server-Side Request Forgery (SSRF) & 22 \\
12 & Allocation of Resources Without Limits or Throttling & 25 \\
\bottomrule
\end{tabular}
\end{center}

\subsubsection{Control de acceso inapropiado}
\begin{practbox}{Idea}
\textbf{Autenticación, autorización y registro.} Aparecen autorizaciones/autenticaciones
incorrectas o \textbf{faltantes}, y formas de \emph{saltear} la autorización. La
\emph{authorization} a secas (\emph{Incorrect Authorization}) ``engloba todo''. Mitigación: el
control de acceso depende mucho de la \textbf{arquitectura de diseño} (dónde poner el control) y
de la \textbf{implementación} (no configurarlo mal); no suele depender de un lenguaje específico.
\end{practbox}
\begin{center}
\renewcommand{\arraystretch}{1.25}
\begin{tabular}{>{\bfseries}c p{8.8cm} c}
\toprule
\# & Debilidad & Rank \\
\midrule
1 & Missing Authorization (autorización faltante) & 4 \\
2 & Improper Access Control (control de acceso inapropiado) & 19 \\
3 & Incorrect Authorization & 17 \\
4 & Missing Authentication for Critical Function & 21 \\
5 & Authorization Bypass Through User-Controlled Key & 24 \\
\bottomrule
\end{tabular}
\end{center}

\subsubsection{Falla en mecanismos de protección}
\begin{practbox}{Idea}
\textbf{Mal uso / no uso de mecanismos de protección.} Incluye el \textbf{uso incorrecto de
algoritmos criptográficos} (o no usar estándares). ``No tenés que inventar tu función de
encripción; tenés que usar los estándares y usarlos correctamente.''
\end{practbox}
\begin{center}
\renewcommand{\arraystretch}{1.25}
\begin{tabular}{>{\bfseries}c p{8.8cm} c}
\toprule
\# & Debilidad & Rank \\
\midrule
1 & Cross-Site Request Forgery (CSRF) & 3 \\
\bottomrule
\end{tabular}
\end{center}

\subsubsection{Otros (cálculos, condiciones excepcionales, flujo de control)}
\begin{center}
\renewcommand{\arraystretch}{1.25}
\begin{tabular}{>{\bfseries}p{4.6cm} >{\bfseries}p{5.5cm} c}
\toprule
Pilar & Debilidad & Rank \\
\midrule
Cálculos incorrectos & Integer Overflow or Wraparound & 30 \\
Condiciones excepcionales (estándares de código) & NULL Pointer Dereference & 13 \\
Manejo insuficiente del flujo de control & Race Condition & 36 \\
\bottomrule
\end{tabular}
\end{center}
Estos aparecen en el \emph{on-the-cusp list} (puestos $\sim$26--36): debilidades que ``están ahí,
queriendo entrar y salir''. Cada una indica su posición del año anterior y cuánto subió/bajó.
\emph{(Ref:} \href{https://cwe.mitre.org/top25/archive/2025/2025_onthecusp_list.html}{\texttt{cwe.mitre.org/top25/archive/2025/2025\_onthecusp\_list.html}}\emph{.)}

\subsection{Debilidades recurrentes (entran y salen del ranking)}
\begin{warnbox}{``Siempre vuelven''}
Algunas debilidades desaparecen un tiempo y \textbf{vuelven a aparecer}:
\begin{itemize}[leftmargin=1.4em,itemsep=1pt]
  \item \textbf{Credenciales hardcoded} (codificadas en el código): hasta el año pasado siempre
        estaban en el ranking.
  \item \textbf{Exposure of Sensitive Information} (exponer información sensible): p.\ ej.\ un
        mensaje de error que filtra parte del directorio o de la base de datos. ``Había
        desaparecido y volvió.''
  \item \textbf{Uso incorrecto / ausencia de criptografía estándar}: estuvo entre los primeros en
        su momento; ``menos mal que ya no aparece''.
\end{itemize}
\textbf{Clásicos que nunca desaparecen}: \emph{Cross-site Scripting} (rank 1, también nº1 el año
pasado) y \emph{SQL Injection} (rank 2, estaba en el 3 el año anterior); fluctúan en los primeros
puestos. La aparición masiva de \textbf{buffer overflows} este año (4--5 entradas) llamó la
atención: a más uso de un lenguaje/registro de datos, más información $\Rightarrow$ más detección.
\end{warnbox}

\subsection{Ejemplo: impacto y fase de un buffer overflow}
\begin{practbox}{Lectura de una ficha CWE (caso de la clase)}
Para \textbf{Buffer Overflow}, la ficha muestra:
\begin{itemize}[leftmargin=1.4em,itemsep=1pt]
  \item \textbf{Impacto (tríada CIA)}: \emph{Integridad} (puede ejecutar comandos no autorizados
        y modificar datos), \emph{Disponibilidad} (puede producir \emph{DoS} / crash) y
        \emph{Confidencialidad} (si se leen datos confidenciales).
  \item \textbf{Fase donde se origina}: \emph{requisitos} (elegir lenguaje/compilador adecuado),
        \emph{arquitectura y diseño} (qué librerías/frameworks usar), \emph{implementación}.
  \item \textbf{Probabilidad de explotación}: alta. Más abajo: relaciones con otros errores,
        lenguajes frecuentes, ejemplos y casos reales observados, y métodos de detección.
\end{itemize}
\end{practbox}

\begin{slidebox}
\textbf{Guía de ejercicios.} La práctica selecciona ejercicios de los tipos más comunes —XSS,
\emph{Race Condition}, \emph{Web sessions}— pero conviene leer todas las fichas \textbf{agrupadas
por categoría} (es más fácil de entender y recordar). Algunos nombres de ataques de canal lateral
(p.\ ej.\ \emph{RAMBO}: variaciones electromagnéticas que forman radiofrecuencias captables desde
afuera) también aparecen en la guía, conectando con la parte teórica.
\end{slidebox}

%=====================================================================
\section*{Lectura recomendada}
\addcontentsline{toc}{section}{Lectura recomendada}
%=====================================================================
\begin{center}
\begin{tcolorbox}[enhanced,colback=itbaBlue!20,colframe=itbaCyanDark,boxrule=1pt,
  arc=3pt,width=0.85\textwidth,halign=center]
\large\textbf{Capítulos 16-1 y 17}\\[4pt]
\emph{Computer Security: Art and Science}\\[2pt]
Matt Bishop
\end{tcolorbox}
\end{center}
\noindent\textbf{Recursos de la práctica:} CWE Top 25 de MITRE
(\href{http://cwe.mitre.org/top25}{\texttt{cwe.mitre.org/top25}}), su jerarquía de \emph{Research
Concepts} (\href{https://cwe.mitre.org/data/definitions/1000.html}{\texttt{cwe.mitre.org/data/definitions/1000.html}})
y OWASP como referencia complementaria de seguridad en aplicaciones web.

%---------------------------------------------------------------------
\vfill
\begin{center}\small\itshape\color{black!60}
Apunte integral de la Clase 9 — Flujo de Información y Malware · Criptografía y Seguridad, ITBA.\\
Síntesis de teoría (transcripciones + diapositivas) y práctica (CWE Top 25).
\end{center}

\end{document}
